\chapter{Appendix: The Real Numbers and Completeness}
\label{ch41:real-numbers-completeness}

Every theorem in this book that begins with the statement ``there exists'' --- a
limit exists, a continuous function attains a maximum, a curve crosses zero,
a bounded monotone sequence converges --- rests on the structure of the
real numbers. Without the completeness of the real number line, these existence
theorems fail, and key parts of calculus break down.

In the previous appendix, we traced the number systems from $\N$ to
$\C$, noting that the passage from fractions to the reals eliminates gaps in the
line. Throughout calculus, we relied on this property when defining
limits, proving that continuous curves cross zero, and optimizing
functions. This appendix examines that foundation directly: \emph{what are the real numbers,
and what makes the real line complete?} We routinely draw the number line
as an unbroken continuum. But why should it have no gaps? The rationals
--- fractions $p/q$ --- already appear to fill the line densely: between any two of them
sits another. Yet the rational numbers contain genuine gaps,
and one of those gaps sits exactly where $\sqrt{2}$ should be.

The property that resolves these gaps is the \textbf{completeness axiom}, also called
the least-upper-bound property. It provides the foundation for the Intermediate Value
Theorem, the Extreme Value Theorem, the convergence of monotone sequences, the existence
of the Riemann integral, and the definition of limits. What follows is the
development of this property and its consequences throughout analysis.

\chapterline

\section{A gap in the rationals}
\label{ch41:gap}

Start with the picture, where the trouble first becomes visible. Draw the rational points
on a line. Between $0$ and $1$ you can mark
$\tfrac12$, and then $\tfrac13, \tfrac23$, and $\tfrac14,\tfrac34$, and on and on. The
marks pile up so thickly that no matter how powerful a magnifying glass you bring, you
will always find infinitely many rationals in any window, however small.

\begin{intuitionbox}
The rationals are \textbf{dense}: between any two distinct rationals $a<b$ lies another
rational, for example their average $\frac{a+b}{2}$. Repeat the averaging and you get
infinitely many. Visually, the rational marks leave no \emph{visible} space anywhere on
the line. It is enormously tempting to conclude that they leave no space at all --- that
the rationals already \emph{are} the continuous line. This is the trap. Density is not
the same as having no holes. A sieve can be dense with metal and still be mostly holes.
\end{intuitionbox}

To demonstrate a gap in the rationals, consider the length the Greeks investigated:
the diagonal of a unit square. By the Pythagorean theorem, a square one unit on a side
has a diagonal of length $d$ with $d^2 = 1^2 + 1^2 = 2$. So the quantity $\sqrt{2}$ is a
definite geometric length constructible with straightedge and compass. The question is
whether that length corresponds to a \emph{rational} number. We prove below that no such rational number exists.

\begin{theorem}[Irrationality of $\sqrt 2$]
\label{ch41:sqrt2-irrational}
There is no rational number $r$ with $r^2 = 2$.
\end{theorem}

\begin{proof}
Suppose, for contradiction, that such a rational exists. Write it in lowest terms as
$r = p/q$, where $p$ and $q$ are integers with no common factor (we can always cancel until
this holds). Then
\[
  \frac{p^2}{q^2} = 2 \quad\Longrightarrow\quad p^2 = 2q^2.
\]
The right-hand side is even, so $p^2$ is even. But the square of an odd number is odd
(an odd number is $2k+1$, and $(2k+1)^2 = 4k^2+4k+1$ is odd), so $p$ itself must be even.
Write $p = 2m$. Substituting,
\[
  (2m)^2 = 2q^2 \quad\Longrightarrow\quad 4m^2 = 2q^2 \quad\Longrightarrow\quad q^2 = 2m^2.
\]
Now the same argument applies to $q$: $q^2$ is even, so $q$ is even. But then $p$ and $q$
are \emph{both} even --- they share the factor $2$ --- contradicting our assumption that
the fraction was in lowest terms. The assumption was impossible. Hence no rational squares
to $2$.
\end{proof}

\begin{connectionbox}
This argument, or something close to it, is among the oldest theorems in mathematics,
attributed to the Pythagoreans of the fifth century BCE. The Pythagorean creed held that
\emph{all is number}, by which they meant whole numbers and their ratios. The discovery
that the diagonal of a square is \emph{incommensurable} with its side --- meaning no common
measuring unit, however tiny, can divide both sides evenly into whole numbers --- was
a significant conceptual crisis. Legend holds that the mathematician who
revealed it, Hippasus, was drowned at sea for revealing the secret. The lasting
realization was that geometry includes lengths that rational fractions cannot express. Closing
that gap between geometry and arithmetic took more than two millennia and was not
completed until the 1870s, with the work of Dedekind and Cantor discussed in the
history box in Section~\ref{ch41:axiom}.
\end{connectionbox}

So there is a definite, constructible geometric point --- the length $\sqrt 2$, roughly
$1.41421356\ldots$ --- that no fraction reaches. The rationals, dense as they are, leave a
gap at that position. By scaling and translating, similar gaps appear everywhere along the line.
The rational line is dense, yet incomplete.

\begin{happeningbox}
Here is the precise structure of the gap, which the completeness axiom will
address. Consider the two sets of rationals
\[
  A = \set{x \in \Q : x > 0 \text{ and } x^2 < 2}, \qquad
  B = \set{x \in \Q : x > 0 \text{ and } x^2 > 2}.
\]
Every positive rational lands in exactly one of $A$ or $B$ (none has square equal to $2$,
by Theorem~\ref{ch41:sqrt2-irrational}). Every element of $A$ is less than every element
of $B$. The set $A$ has no \emph{largest} member: given any $x\in A$, one can always choose
a larger rational still below $\sqrt 2$. The set $B$ has no \emph{smallest} member: any
$x\in B$ can be nudged down slightly and stay above $\sqrt 2$. So the rational line
partitions into a left set with no maximum and a right set with no minimum, with
\emph{no rational number sitting between them}. This partition of the ordered numbers into two
complementary sets is what Richard Dedekind called a \textbf{cut} (Figure~\ref{ch41:fig-dedekind-cut}).
In $\R$, there is always a number located at such a boundary; in $\Q$, there is sometimes none.
That distinction is completeness.
\end{happeningbox}

\begin{figure}[H]
\centering
\begin{tikzpicture}[scale=1.0]
  % Number line axis
  \draw[->,thick] (-0.5,0) -- (10.5,0) node[right] {$\Q$};

  % Split point at x = 5.0 (representing sqrt(2))
  \def\cutx{5.0}

  % Left set A (x^2 < 2)
  \draw[sourceblue, line width=3pt] (0.2,0) -- (\cutx,0);
  \fill[sourceblue] (0.2,0) circle (2pt) node[below=3pt] {$0$};
  \fill[white,draw=sourceblue,thick] (\cutx,0) circle (3pt);

  % Right set B (x^2 > 2)
  \draw[targetgreen!80!black, line width=3pt] (\cutx,0) -- (9.8,0);
  \fill[white,draw=targetgreen!80!black,thick] (\cutx,0) circle (3pt);

  % Scissors / cut line
  \draw[trapred, dashed, thick] (\cutx, -1.1) -- (\cutx, 1.4);
  \fill[trapred] (\cutx, 0) circle (2.2pt);
  \node[trapred, above=2pt, font=\small] at (\cutx, 1.4) {Cut at $\sqrt{2}$ (not in $\Q$)};

  % Brackets / labels for Set A and Set B
  \node[sourceblue, font=\small, align=center] at (2.5, 0.9)
    {Set $A = \{x \in \Q : x^2 < 2\}$\\(\emph{No maximum element})};
  \node[targetgreen!80!black, font=\small, align=center] at (7.6, 0.9)
    {Set $B = \{x \in \Q : x^2 > 2\}$\\(\emph{No minimum element})};

  % Rational tick samples
  \foreach \x/\lab in {1.8/{1}, 3.6/{1.4}, 4.4/{1.41}, 5.6/{1.42}, 6.4/{1.5}, 8.2/{2}} {
    \draw (\x, 0.08) -- (\x, -0.08) node[below=3pt, font=\footnotesize] {$\lab$};
  }

  % Annotation of empty seam
  \node[trapred, font=\footnotesize, align=center] at (\cutx, -1.6)
    {Empty seam in $\Q$: nothing sits between $A$ and $B$.\\Completeness inserts the real number $\sqrt{2}$ here.};
\end{tikzpicture}
\caption{A Dedekind cut dividing the rational numbers at $\sqrt{2}$. Set $A$ contains all positive rationals with $x^2 < 2$ (blue); set $B$ contains all rationals with $x^2 > 2$ (green). Because $\sqrt{2}$ is irrational, set $A$ has no largest element and set $B$ has no smallest element. The seam between them is completely empty in $\Q$. Dedekind's breakthrough was defining the real number $\sqrt{2}$ to \emph{be} this partition.}
\label{ch41:fig-dedekind-cut}
\end{figure}

\subsection*{Squeezing the hole between fractions}

We can watch the hole numerically. Build two sequences of rationals, one climbing toward
$\sqrt 2$ from below and one descending toward it from above, by choosing more and more
decimal digits:
\[
\begin{array}{c|c|c}
n & \text{below: } a_n & \text{above: } b_n \\ \hline
1 & 1.4 & 1.5 \\
2 & 1.41 & 1.42 \\
3 & 1.414 & 1.415 \\
4 & 1.4142 & 1.4143 \\
5 & 1.41421 & 1.41422
\end{array}
\]
Each $a_n$ is a rational with $a_n^2 < 2$ and each $b_n$ is a rational with $b_n^2 > 2$,
and the gap $b_n - a_n = 10^{-n}$ shrinks to zero. The two sequences close in on a single
location on the line --- you can pin it to within any tolerance you like --- and yet
\emph{that location is occupied by no rational at all}. The fractions march right up to the
edge of a point that, in their world, simply is not there.

\begin{figure}[H]
\centering
\begin{tikzpicture}
  % number line
  \draw[->,thick] (-0.5,0) -- (10.5,0);
  \foreach \x/\lab in {0/1.40, 2.5/1.41, 5/1.42, 7.5/1.43, 10/1.44}
    \draw (\x,0.12) -- (\x,-0.12) node[below] {\small \lab};
  % the true point sqrt2 ~ 1.41421, map 1.40->0, 1.44->10 => scale 250 per 0.01
  % position of 1.41421 : (1.41421-1.40)*250 = 0.01421*250 = 3.5525
  \def\px{3.5525}
  \fill[trapred] (\px,0) circle (2.3pt);
  \node[trapred, above] at (\px,0.55) {\small $\sqrt 2$ (no fraction here)};
      \draw[trapred, ->, thin] (\px,0.55) -- (\px,0.10);
  % a_n marks climbing from left
  \foreach \x in {2.5, 3.0, 3.35, 3.50}
    \draw[sourceblue, thick] (\x,0.05) -- (\x,0.30);
  % b_n marks descending from right
  \foreach \x in {5.0, 4.0, 3.75, 3.60}
    \draw[targetgreen, thick] (\x,-0.05) -- (\x,-0.30);
  \node[sourceblue, font=\small] at (1.1,0.45) {rationals $a_n < \sqrt{2}$};
  \node[targetgreen, font=\small] at (6.2,-0.60) {rationals $b_n > \sqrt{2}$};
\end{tikzpicture}
\caption{Rationals from the left ($a_n$, blue) and from the right ($b_n$, green) close in
on the location of $\sqrt 2$. The gap $b_n-a_n=10^{-n}$ shrinks to zero, pinning a single
point --- but in the rational line that point is a hole. Completeness is the axiom that
puts a real number there.}
\label{ch41:fig-squeeze}
\end{figure}

\begin{whybox}
Why does a missing point matter when rationals suffice for any physical
measurement? Because calculus relies on general theorems guaranteeing that limits
exist before we can compute them in closed form. The rational sequence
$1.4, 1.41, 1.414, \ldots$ is strictly increasing and bounded above by $2$, yet in
$\Q$ it has no limit at all. To guarantee that bounded monotone sequences and
Cauchy sequences actually converge, the number system cannot have missing targets.
\end{whybox}

\section{Upper bounds, suprema, and the language of ceilings}
\label{ch41:sup}

To state the axiom that fixes everything, we first need vocabulary for ``the place a set
is heading toward from above.'' This is the notion of a \emph{least upper bound}, and it
deserves to be built slowly: it is subtle, and it underlies everything that follows.

\begin{definition}[Bounded above; upper bound]
\label{ch41:def-bounded}
Let $S$ be a nonempty set of real numbers. A number $M$ is an \textbf{upper bound} for
$S$ if $x \le M$ for every $x \in S$. If $S$ has at least one upper bound, we say $S$ is
\textbf{bounded above}. \emph{Lower bound} and \textbf{bounded below} are defined the same
way with the inequality reversed.
\end{definition}

An upper bound is a ceiling: a number that no element of $S$ exceeds. A set can have many
ceilings. The interval $S = (0,1)$ has $1$ as a ceiling, but also $2$, also $17.5$, also a
billion. Ceilings come in a whole upward family; once one number is a ceiling, so is every
number above it.

The interesting ceiling is the \emph{lowest} one --- the tightest, the one that hugs the
set from above with no slack.

\begin{definition}[Supremum / least upper bound]
\label{ch41:def-sup}
A number $L$ is the \textbf{supremum} (or \textbf{least upper bound}) of a nonempty set
$S$, written $L = \sup S$, if both:
\begin{enumerate}
  \item $L$ is an upper bound for $S$ (nothing in $S$ exceeds $L$), and
  \item $L$ is the \emph{least} such bound: if $M$ is any upper bound for $S$, then
        $L \le M$.
\end{enumerate}
The mirror notion, the \textbf{infimum} or greatest lower bound, is written $\inf S$.
It is the highest possible floor that stays beneath every point in the set.
\end{definition}

\begin{intuitionbox}
Think of lowering a ceiling slowly toward a crowd of people. You keep lowering until,
just before it would touch someone's head, you stop. That stopping height is the
supremum: the lowest you can put the ceiling while still clearing everyone. Two flavours
can happen at the moment you stop. Either the ceiling is resting on the tallest person's
head (the supremum belongs to the set --- it is a \emph{maximum}), or there is no single
tallest person, only an unending sequence of ever-taller people whose heights approach
the ceiling without any one reaching it (the supremum is the limit of their heights but
belongs to no individual). The supremum captures \emph{both} cases in one word, which is
exactly why it is more powerful than ``maximum.''
\end{intuitionbox}

The second flavour is the one to internalise, because it is where the supremum earns its
keep. Consider the open interval $S = (0,1) = \set{x : 0 < x < 1}$. There is no largest
element: whatever $x<1$ you name, the number $\frac{x+1}{2}$ is larger and still below
$1$. So $S$ has no maximum. But it plainly has a least upper bound: $\sup S = 1$. The
number $1$ is an upper bound, and nothing smaller can be, because any $M<1$ is itself
beaten by an element of $S$ (the point $\frac{M+1}{2}$ lies in $S$ and exceeds $M$). The
supremum exists even though the maximum does not. \emph{This is precisely why the
concept was invented.}

\begin{figure}[H]
\centering
\begin{tikzpicture}
  \draw[->,thick] (-0.5,0) -- (8.5,0);
  \draw (0,0.1) -- (0,-0.1) node[below] {\small $0$};
  % open interval (0,1)
  \draw[sourceblue, very thick] (0,0) -- (6,0);
  \draw[sourceblue, fill=white, thick] (0,0) circle (2.5pt);
  \draw[sourceblue, fill=white, thick] (6,0) circle (2.5pt);
  \node[sourceblue] at (3,0.5) {$S=(0,1)$};
  % a few upper bounds strictly above the supremum
  \foreach \x in {6.8, 7.6}
    \draw[quotientorange, thick, ->] (\x,0.7) -- (\x,0.15);
  \node[quotientorange, font=\small, align=center] at (7.2,1.05) {upper\\bounds};
  % sup marker
  \draw[targetgreen, very thick] (6,-0.55) -- (6,0.55);
  \node[targetgreen, below] at (6.0,-0.65) {\small $\sup S = 1$ (not in $S$)};
\end{tikzpicture}
\caption{The open interval $(0,1)$ has no largest element --- the right endpoint is
missing (hollow circle) --- yet it has a least upper bound, $\sup S = 1$. The orange
arrows mark some of the upper bounds; the supremum is the leftmost of them, sitting
exactly at the missing edge.}
\label{ch41:fig-sup}
\end{figure}

\begin{focusbox}
A clean way to test a candidate $L$ for being $\sup S$ uses two complementary promises,
and you will use this constantly in proofs:
\begin{enumerate}
  \item \textbf{Ceiling:} $x \le L$ for all $x\in S$. (Nothing pokes above $L$.)
  \item \textbf{Tightness:} for every $\eps>0$ there is some $x\in S$ with $x > L-\eps$.
        (You cannot lower the ceiling by any amount $\eps$, however tiny, without
        cutting into the set.)
\end{enumerate}
The second promise is the operational heart of the supremum. It says: elements of $S$ come
\emph{arbitrarily close} to $L$ from below. Whenever you see ``$\sup$'' in a proof,
reach for this $\eps$ characterisation --- it converts the abstract ``least upper bound''
into a concrete sequence of elements creeping up to $L$, which is something you can
compute with.
\end{focusbox}

\begin{recoverybox}
\textbf{Pause and think:} Consider the sequence set $S = \set{2 - \frac{1}{n} : n \in \N} = \set{1, \frac{3}{2}, \frac{5}{3}, \frac{7}{4}, \ldots}$.
What is $\inf S$, and is it attained? What is $\sup S$, and is it attained?
\emph{(Check: The smallest term is the first one, so $\inf S = 1$, which is attained. As $n$ grows, the terms climb toward $2$ but never reach it. Lowering the ceiling to $2 - \eps$ always leaves terms $2 - 1/n$ above it once $n > 1/\eps$, so $\sup S = 2$, which is not attained.)}
\end{recoverybox}

\begin{figure}[H]
\centering
\begin{tikzpicture}[scale=1.05]
  % Axis
  \draw[->,thick] (0,0) -- (10.5,0) node[right] {$\R$};

  % Sup coordinate at 8.0
  \def\supx{8.0}
  \def\epsval{2.4}
  \def\lowx{5.6} % supx - epsval

  % Epsilon trap zone [L - eps, L] drawn first
  \fill[quotientorange!20, rounded corners=2pt] (\lowx, -0.55) rectangle (\supx, 0.55);
  \draw[quotientorange, thick, dashed] (\lowx, -0.75) -- (\lowx, 0.75) node[above, font=\small] {$L - \eps$};

  % Sample set points
  \foreach \x in {1.2, 2.1, 3.4, 4.2, 5.0} {
    \fill[sourceblue] (\x, 0) circle (2.2pt);
  }
  \node[sourceblue, font=\small] at (3.0, 0.45) {Elements of set $S$};

  % Points trapped in the epsilon band
  \fill[trapred] (6.2, 0) circle (2.6pt);
  \fill[trapred] (7.1, 0) circle (2.6pt) node[above=5pt, font=\small\bfseries, trapred!80!black] {$x \in S$};
  \fill[trapred] (7.7, 0) circle (2.6pt);

  % Ceiling line at L = sup S
  \draw[targetgreen!80!black, line width=1.5pt] (\supx, -0.9) -- (\supx, 1.2);
  \node[targetgreen!80!black, font=\small\bfseries, above] at (\supx, 1.2) {Ceiling: $L = \sup S$};
  \draw[targetgreen!80!black, fill=white, thick] (\supx, 0) circle (3pt);

  % Epsilon bracket
  \draw[decorate, decoration={brace, amplitude=5pt}, quotientorange, thick]
    (\supx, -0.7) -- (\lowx, -0.7) node[midway, below=5pt, font=\small] {Tolerance zone of width $\eps$};

  % Annotation
  \node[anchor=west, font=\small, trapred!80!black] at (0.5, -1.35)
    {\textbf{Tightness Promise:} Lowering the ceiling to $L - \eps$ \emph{always} traps at least one element $x \in S$.};
\end{tikzpicture}
\caption{The $\eps$-characterisation of the supremum. $L = \sup S$ acts as an impassable ceiling for $S$. Tightness asserts that no matter how tiny the tolerance $\eps > 0$ is, lowering the ceiling to $L - \eps$ breaches the set: there must exist an element $x \in S$ inside the orange corridor $(L - \eps, L]$.}
\label{ch41:fig-sup-eps-trap}
\end{figure}

\begin{trapbox}
Do not confuse the supremum with the maximum, and do not assume the supremum belongs to
the set. For $(0,1)$, the supremum is $1$ but $1\notin(0,1)$. For $[0,1]$, the supremum is
$1$ and $1\in[0,1]$, so here it is also the maximum. A maximum is a supremum that happens
to be attained; a supremum need not be attained. The reverse trap also bites: a set can
have a supremum without a maximum, but it can never have a maximum without that maximum
being the supremum. ``Max exists'' is the strong statement; ``sup exists'' is the one the
axiom will guarantee.
\end{trapbox}

\section{The completeness axiom}
\label{ch41:axiom}

We have seen exactly what goes wrong in $\Q$: the set $A = \set{x\in\Q : x>0,\ x^2<2}$ is
bounded above (by $2$, say) but has no least upper bound \emph{inside $\Q$} --- its only
conceivable least upper bound is $\sqrt 2$, which $\Q$ does not contain. The seam between
$A$ and the larger rationals is empty. The real numbers are, by definition, the system in
which this never happens.

\begin{focusbox}
\textbf{The Completeness Axiom (least-upper-bound property).}
\emph{Every nonempty set of real numbers that is bounded above has a least upper bound
in $\R$.} In symbols: if $S\subseteq\R$ is nonempty and bounded above, then $\sup S$
exists and is a real number.
\end{focusbox}

That is the entire axiom. It does not mention $\sqrt 2$, limits, or
continuity; it speaks only of sets and upper bounds. Yet it is the defining property that
distinguishes $\R$ from $\Q$, and it underlies every existence theorem in
calculus.

\begin{happeningbox}
What does the axiom mean geometrically? It asserts that \textbf{the real line has no gaps}.
Recall the cut from Section~\ref{ch41:gap}: a left subset $A$ with no maximum facing a right subset
with no minimum, and no rational number between them. Completeness ensures that this cannot happen
in $\R$. If $S$ is bounded above, the upper bounds of $S$ form a set to the right, and the elements
of $S$ lie to the left; completeness requires that the boundary between elements of $S$ and upper
bounds of $S$ is \emph{an actual real number}, $\sup S$. Every cut determines a unique real number.
This is why the real line forms an unbroken continuum: completeness is the precise mathematical
formulation of that property.
\end{happeningbox}

\begin{remark}
By symmetry (replace $S$ with $-S = \set{-x : x\in S}$ and flip every inequality), the
axiom has a twin: \emph{every nonempty set bounded below has a greatest lower bound,}
$\inf S$, in $\R$. Suprema and infima always exist for bounded sets. You only ever need to
prove the version you find convenient; the other follows for free by reflection.
\end{remark}

\begin{historybox}
Why is completeness an \emph{axiom} and not a theorem? Because whether you take it as an
axiom or prove it as a theorem depends on where you choose to start.

The trouble began in the fifth century BCE when the Pythagoreans discovered that the diagonal of
a unit square is \emph{incommensurable} with its side --- no common unit of measurement, however
fine, can divide both sides into whole numbers. For two thousand years, geometry had to stand in
for numbers because arithmetic had no language for irrational quantities. When Newton and Leibniz
invented calculus in the late seventeenth century, they relied on geometric intuition of continuous
curves, fluent quantities, and infinitesimals without any formal arithmetic definition of the continuum.
In 1821, Augustin-Louis Cauchy published his \emph{Cours d'Analyse}, attempting to place calculus on
a solid foundation of limits. Yet Cauchy tacitly assumed without proof that whenever the terms of a
sequence crowd closer and closer to one another (what we now call a \textbf{Cauchy sequence}), they
must inevitably converge to an actual number.

The crisis came to a head in the autumn of 1858. Richard Dedekind, then preparing lecture notes
on differential calculus at Zurich Polytechnic, was horrified to realize that he had no rigorous
way to prove that an increasing bounded sequence has a limit without appealing to the physical,
geometric drawing of a line. But using geometry to justify arithmetic is circular: analytic
geometry uses coordinate pairs $(x,y) \in \R^2$ to define the line in the first place! Arithmetic had to
stand on its own feet. Armed with only the rational numbers $\Q$ and basic set operations,
mathematicians set out to construct the continuum from scratch.

In the 1870s, four mathematicians independently succeeded: Richard Dedekind, Georg Cantor,
Charles M\'eray, and Karl Weierstrass. Dedekind (published in 1872 in \emph{Stetigkeit und
irrationale Zahlen}) defined a real number to be a ``cut''---a partition of $\Q$ into a lower set
and an upper set, placing an actual number into every empty seam. Cantor instead defined a real number
as an equivalence class of Cauchy sequences of rationals, treating sequences that close in on the
same spot as the same number. If you define the reals by Dedekind cuts or Cauchy sequences,
the least-upper-bound property becomes a provable theorem. If you start with $\R$ as an ordered
field, you take it as an axiom.

Yet this breakthrough sparked ferocious controversy. Leopold Kronecker, a strict finitary
constructivist---one who believed that mathematical objects exist only if they can be calculated
in a finite number of algorithmic steps---waged an intellectual war against Cantor and Dedekind.
Kronecker famously proclaimed: \emph{``Die ganzen Zahlen hat der liebe Gott gemacht, alles andere
ist Menschenwerk''} (``God made the integers, all else is the work of man''). Kronecker argued
that an irrational number like $\sqrt{2}$ has no legitimate existence unless accompanied by a
finite algorithm to compute it, rejected completed infinities, and attempted to suppress Cantor's
work.

History confirmed the value of Dedekind and Cantor's work: without their rigorous model of the
continuum, modern analysis, topology, and differential equations could not have developed
their current foundations. Yet foundational questions remain active. The philosophical debate
between classical mathematics (which accepts non-constructive existence by completeness) and
constructive mathematics (which requires explicit algorithms) continues in logic and theoretical
computer science. In addition, Cantor's discovery that the reals are uncountably infinite led to
the \textbf{Continuum Hypothesis}---the question of whether there exists an infinite set strictly
larger than $\N$ but strictly smaller than $\R$. In 1940 Kurt G\"odel and in 1963 Paul Cohen
proved that the Continuum Hypothesis is formally undecidable from the standard axioms of set theory (ZFC),
demonstrating that the question cannot be settled within the standard ZFC framework alone.
\end{historybox}

\subsection*{Completeness builds $\sqrt 2$}

To see how the axiom operates, we use it to prove the existence of $\sqrt 2$.

\begin{theorem}[Existence of square roots]
\label{ch41:sqrt-exists}
There is a positive real number $L$ with $L^2 = 2$.
\end{theorem}

\begin{proof}[Proof sketch]
Let $S = \set{x\in\R : x > 0 \text{ and } x^2 < 2}$. The set is nonempty (it contains $1$,
since $1^2=1<2$) and bounded above (by $2$, since any $x>2$ has $x^2>4>2$ and so cannot be
in $S$). By the completeness axiom, $L=\sup S$ exists as a real number. We claim $L^2=2$.

Suppose instead $L^2 < 2$. Then $L$ is too small to be the ceiling: we can nudge it up and
stay inside $S$. Concretely, for small $h>0$,
\[
  (L+h)^2 = L^2 + 2Lh + h^2 \le L^2 + (2L+1)h \quad\text{(taking $h\le 1$),}
\]
and choosing $h$ small enough that $(2L+1)h < 2 - L^2$ makes $(L+h)^2 < 2$, so
$L+h\in S$ --- but $L+h>L$ contradicts $L$ being an \emph{upper} bound.

Suppose instead $L^2 > 2$. Then $L$ is too big to be the \emph{least} ceiling: we can
shave it down and still have a ceiling. For small $h>0$ (taking $h < L$ so $L-h > 0$),
\[
  (L-h)^2 = L^2 - 2Lh + h^2 > L^2 - 2Lh,
\]
and choosing $h$ small enough that $2Lh < L^2 - 2$ ensures $(L-h)^2 > 2$. For any $x\in S$,
we have $x^2 < 2 < (L-h)^2$. Because both $x$ and $L-h$ are positive, this implies $x < L-h$.
Thus $L-h$ is an upper bound for $S$ --- but $L-h<L$ contradicts $L$ being the \emph{least}
upper bound.

Both inequalities are impossible, so $L^2 = 2$. The completeness axiom has produced the
number the rationals were missing.
\end{proof}

\begin{aestheticbox}
The proof above illustrates the power of the completeness axiom. We did not
construct $\sqrt 2$ digit by digit. Instead, we defined a set --- the positive
numbers whose square is less than $2$ --- and the axiom guaranteed that its least
upper bound exists in $\R$, with that bound satisfying $L^2 = 2$.
The supremum provides a general existence principle: for any nonempty bounded set,
an exact supremum is guaranteed. Most existence theorems in calculus
follow this pattern. By defining appropriate sets --- points where a function is
negative, function values on a closed interval, or partial sums of a series ---
completeness guarantees the existence of the critical value: the crossing point,
the maximum, the limit, or the integral.
\end{aestheticbox}

\begin{recoverybox}
\textbf{Pause and think:} Does the set of all negative real numbers $S = (-\infty, 0) = \set{x \in \R : x < 0}$ have a supremum? Is it bounded above? Does the completeness axiom apply?
\emph{(Check: Any number $\ge 0$ is an upper bound, so $S$ is bounded above, and $S$ is nonempty. Completeness guarantees that $\sup S$ exists. The least upper bound is $\sup S = 0$, which is not in $S$.)}
\end{recoverybox}

\section{Why calculus needs completeness}
\label{ch41:needs}

Four fundamental theorems in calculus follow directly from the completeness axiom.
We prove the first two in full and outline the others, showing how completeness
supplies the required existence guarantees in each case.

\subsection*{Monotone sequences converge}

\begin{theorem}[Monotone Convergence Theorem]
\label{ch41:mct}
A sequence $(a_n)$ that is increasing ($a_1 \le a_2 \le a_3 \le \cdots$) and bounded above
converges to a real limit. In fact $a_n \to \sup\set{a_n : n\in\N}$.
\end{theorem}

\begin{proof}
Let $S = \set{a_n : n\in\N}$, the set of values the sequence takes. It is nonempty and
bounded above by hypothesis, so by completeness $L=\sup S$ exists. We show $a_n\to L$.

Fix any tolerance $\eps>0$. By the tightness property of the supremum
(Section~\ref{ch41:sup}), $L-\eps$ is not an upper bound, so some term $a_N$ satisfies
$a_N > L-\eps$. Because the sequence is increasing, every later term is at least as large:
for all $n\ge N$,
\[
  L-\eps < a_N \le a_n \le L < L + \eps,
\]
the last inequality because $L$ is an upper bound for all the $a_n$. Hence
$\abs{a_n - L}<\eps$ for all $n\ge N$. Since $\eps$ was arbitrary, $a_n\to L$
(Figure~\ref{ch41:fig-mct}).
\end{proof}

\begin{figure}[H]
\centering
\begin{tikzpicture}
\begin{axis}[
  width=11cm, height=6cm,
  axis lines=middle,
  xlabel={$n$}, ylabel={$a_n$},
  xlabel style={at={(1,0)}, anchor=south west},
  ylabel style={at={(ticklabel* cs:1)}, anchor=south west},
  xmin=0, xmax=9.5, ymin=-0.15, ymax=1.25,
  xtick={1,2,3,4,5,6,7,8}, xticklabels={$1$,$2$,$3$,$4$,$5$,$6$,$7$,$N{=}8$}, ytick={0,0.5,1},
  yticklabels={$0$,$0.5$,$L=\sup$},
  clip=false,
]
  % the ceiling L = sup
  \addplot[targetgreen, thick, dashed, domain=0:9.5] {1};
  \node[targetgreen] at (axis cs:8.6,1.08) {\small $L$};
  % the band L-eps
  \addplot[quotientorange, thick, dotted, domain=0:9.5] {0.86};
  \node[quotientorange, right] at (axis cs:9.5,0.86) {\small $L-\eps$};
  % the increasing bounded sequence a_n = 1 - 1/n
  \addplot[only marks, sourceblue, mark=*, mark size=1.6pt]
    coordinates {(1,0)(2,0.5)(3,0.6667)(4,0.75)(5,0.8)(6,0.8333)(7,0.8571)(8,0.875)};
  \node[sourceblue] at (axis cs:5.5,0.55) {\small $a_n$ increasing};
  % the first index past the band
  \draw[gray, dashed] (axis cs:8,0.875) -- (axis cs:8,-0.05);
  
\end{axis}
\end{tikzpicture}
\caption{Monotone Convergence in one picture. An increasing sequence $a_n$ (blue dots) that
is bounded above presses upward against its supremum $L$ (green dashed ceiling). For any
tolerance $\eps$, tightness of the supremum forces some term $a_N$ above the line $L-\eps$
(orange dotted); from that index on, every term is trapped in the strip $(L-\eps,L]$, so
$a_n\to L$.}
\label{ch41:fig-mct}
\end{figure}

\begin{connectionbox}
This is the theorem that establishes convergence for the sequence $1.4, 1.41, 1.414,\ldots$ from
Section~\ref{ch41:gap}. It is increasing and bounded above by $2$, so by Monotone
Convergence it converges, and its limit is $\sqrt 2$ as established
in Theorem~\ref{ch41:sqrt-exists}. Decimal expansions rely fundamentally
on this theorem: the digits $3.14159\ldots$ define an increasing,
bounded sequence of rationals, and completeness guarantees that this sequence
converges to a well-defined real number $\pi$. Without completeness, infinite
decimal expansions would not be guaranteed to represent real numbers.
\end{connectionbox}

\subsection*{The Intermediate Value Theorem}

\begin{theorem}[Intermediate Value Theorem]
\label{ch41:ivt}
If $f$ is continuous on $[a,b]$ and $f(a) < 0 < f(b)$, then there is a point
$c\in(a,b)$ with $f(c)=0$.
\end{theorem}

\begin{proof}[Proof via completeness]
Consider the set of inputs on which $f$ is still negative:
\[
  S = \set{x \in [a,b] : f(x) < 0}.
\]
The set is nonempty ($a\in S$, since $f(a)<0$) and bounded above (by $b$). By the completeness
axiom, $c=\sup S$ exists as a real number, and $a \le c \le b$. Because $f(a) < 0$ and $f(b) > 0$,
continuity at the endpoints ensures $c \in (a,b)$. We claim $f(c)=0$, and we rule out both
inequalities using continuity.

If $f(c)<0$: by continuity, choose $\eps = -f(c)/2 > 0$. There exists $\delta>0$ such that $f(x)<0$
for all $x \in (c-\delta, c+\delta) \cap [a,b]$. In particular, the point $c+\delta/2$ satisfies
$f(c+\delta/2)<0$, so $c+\delta/2 \in S$. But $c+\delta/2 > c$, contradicting the fact that $c$ is
an upper bound for $S$.

If $f(c)>0$: by continuity, choose $\eps = f(c)/2 > 0$. There exists $\delta>0$ such that $f(x)>0$
for all $x \in (c-\delta, c+\delta) \cap [a,b]$. Then no point in $(c-\delta, c]$ can belong to $S$.
Because $c$ bounds $S$ from above, every element $x \in S$ must satisfy $x \le c-\delta$.
Thus $c-\delta$ is an upper bound for $S$. But $c-\delta < c$, contradicting the fact that $c$ is
the \emph{least} upper bound.

Both inequalities fail, so $f(c)=0$. The crossing point is the supremum of the ``still negative'' set.
\end{proof}

\begin{figure}[H]
\centering
\begin{tikzpicture}
\begin{axis}[
  width=11cm, height=6.2cm,
  axis lines=middle,
  xlabel={$x$}, ylabel={$f(x)$},
  xtick={1,5}, xticklabels={$a$, $b$},
  ytick=\empty,
  xmin=0.4, xmax=5.6, ymin=-2.2, ymax=2.6,
  samples=120, domain=1:5,
  clip=false,
]
  % a continuous curve going from negative to positive, crossing once
  \addplot[curveblue, very thick] {0.9*(x-3) + 0.45*sin(deg(2.2*(x-3)))};
  % zero crossing near x=3
  \addplot[gray, dashed, domain=0.4:5.6] {0};
  \node[curveblue] at (axis cs:4.6,2.0) {$y=f(x)$};
  % mark f(a)<0 and f(b)>0
  \fill[trapred] (axis cs:1, {0.9*(1-3)+0.45*sin(deg(2.2*(1-3)))}) circle (2.2pt);
  \node[trapred, left] at (axis cs:1,{0.9*(1-3)+0.45*sin(deg(2.2*(1-3)))}) {\small $f(a)<0$};
  \fill[targetgreen] (axis cs:5, {0.9*(5-3)+0.45*sin(deg(2.2*(5-3)))}) circle (2.2pt);
  \node[targetgreen, right] at (axis cs:5,{0.9*(5-3)+0.45*sin(deg(2.2*(5-3)))}) {\small $f(b)>0$};
  % the crossing
  \fill[structurepurple] (axis cs:3,0) circle (2.5pt);
  \node[structurepurple, above left] at (axis cs:3,0) {\small $c=\sup S$};
  % shade S (where f<0) on the axis
  \draw[sourceblue, very thick] (axis cs:1,-0.12) -- (axis cs:3,-0.12);
  \node[sourceblue, below] at (axis cs:2,-0.18) {\small $S=\set{x:f(x)<0}$};
\end{axis}
\end{tikzpicture}
\caption{The Intermediate Value Theorem as a supremum. The curve starts below the axis and
ends above it. The set $S$ where $f$ is still negative (blue segment) is bounded above;
completeness guarantees $c=\sup S$ exists, and continuity forces $f(c)=0$ --- the curve must cross
exactly at the edge of $S$.}
\label{ch41:fig-ivt}
\end{figure}

\begin{whybox}
Why does the IVT \emph{require} completeness and not just continuity? Because the very same
continuous formula can have a crossing or not, depending only on whether the line has a
hole there. Take $f(x)=x^2-2$ on $[0,2]$. It is continuous, $f(0)=-2<0$, and $f(2)=2>0$. As
a function of \emph{real} inputs it crosses zero at $x=\sqrt2$. But as a function of
\emph{rational} inputs --- exactly the same formula, restricted to $\Q$ --- it never equals
zero, because $\sqrt2\notin\Q$. The function jumps from negative to positive
\emph{without ever passing through zero}, because the zero is in a hole. Continuity is
satisfied either way; what changes the conclusion is whether the underlying number system
has the point $\sqrt2$. The IVT is false over $\Q$ and true over $\R$, and the only
difference between $\Q$ and $\R$ is completeness. So completeness is not a technicality
behind the IVT --- it \emph{is} the IVT.
\end{whybox}

\subsection*{The Extreme Value Theorem}

\begin{theorem}[Extreme Value Theorem, sketch]
\label{ch41:evt}
A function continuous on a closed bounded interval $[a,b]$ attains a maximum and a minimum
value on $[a,b]$.
\end{theorem}

The proof proceeds in two steps, using completeness in each. First, one shows that a continuous
function on $[a,b]$ is \emph{bounded}. If it were unbounded, there would exist inputs $x_n$
with $f(x_n)\to\infty$, and completeness (via the Bolzano--Weierstrass theorem, below)
yields a convergent subsequence whose limit produces a contradiction with continuity. Second,
knowing the range $\set{f(x):x\in[a,b]}$ is bounded above, completeness guarantees that
$M=\sup\set{f(x):x\in[a,b]}$ exists in $\R$, and a similar subsequence argument produces a
point $c \in [a,b]$ with $f(c)=M$. The maximum \emph{value} is a supremum (which always exists by
completeness); continuity on a \emph{closed} interval ensures that this supremum is \emph{attained}.
Drop completeness and the bound $M$ might not exist; drop closedness and the supremum might not
be attained, exactly as $\sup(0,1)=1$ is never reached inside $(0,1)$.

This theorem provides the foundation for the Mean Value Theorem in Chapter~15. The proof of the
Mean Value Theorem relies on Rolle's Theorem, which requires the function to attain an interior
extremum. Without the Extreme Value Theorem, Rolle's Theorem cannot be applied; without completeness,
the Extreme Value Theorem does not hold. Order completeness is thus the fundamental property
supporting the Mean Value Theorem.

\begin{theorem}[Bolzano--Weierstrass, sketch]
\label{ch41:bw}
Every bounded sequence of real numbers has a convergent subsequence.
\end{theorem}

The idea is a binary hunt. Enclose the bounded sequence in an interval $[A_0,B_0]$. Cut it
in half; at least one half contains infinitely many terms of the sequence --- keep that
half and call it $[A_1,B_1]$. Halve again, keep an infinite half, and repeat forever. The
left endpoints $A_0\le A_1\le A_2\le\cdots$ form an increasing bounded sequence, so by
Monotone Convergence (Theorem~\ref{ch41:mct}, itself completeness) they converge to a point
$\ell$; the interval lengths halve to zero, so the right endpoints converge to the same
$\ell$. Picking one sequence term from each nested interval yields a subsequence trapped in
intervals shrinking to $\ell$, hence converging to $\ell$. The nested intervals close on a
single point precisely \emph{because the line has no holes} for them to close on
emptily --- completeness again.

\begin{figure}[H]
\centering
\begin{tikzpicture}
  % level 0
  \draw[thick] (0,0) -- (10,0);
  \node[left] at (0,0) {\small $A_0$}; \node[right] at (10,0) {\small $B_0$};
  % level 1 keep right half (5..10)
  \draw[sourceblue, very thick] (5,-0.8) -- (10,-0.8);
  \node[left] at (5,-0.8) {\small $A_1$}; \node[right] at (10,-0.8) {\small $B_1$};
  % level 2 keep left half of that (5..7.5)
  \draw[sourceblue, very thick] (5,-1.6) -- (7.5,-1.6);
  \node[left] at (5,-1.6) {\small $A_2$}; \node[right] at (7.5,-1.6) {\small $B_2$};
  % level 3 keep right half (6.25..7.5)
  \draw[sourceblue, very thick] (6.25,-2.4) -- (7.5,-2.4);
  \node[left] at (6.25,-2.4) {\small $A_3$};
  % limit point
  \draw[trapred, dashed] (6.9,0.3) -- (6.9,-2.7);
  \fill[trapred] (6.9,-2.4) circle (2pt);
  \node[trapred, below] at (6.9,-2.7) {\small $\ell$};
\end{tikzpicture}
\caption{Bolzano--Weierstrass by repeated halving. Each interval keeps a half containing
infinitely many sequence terms; the nested intervals shrink to a single point $\ell$
(dashed line), which exists because completeness leaves no empty seam for them to collapse
into. A subsequence trapped in these intervals converges to $\ell$.}
\label{ch41:fig-bw}
\end{figure}

\subsection*{The integral exists}

When we built the definite integral in Chapter~20 as a limit of Riemann sums, we quietly used
completeness twice. For any partition $P = \{a = x_0 < x_1 < \cdots < x_n = b\}$, the
\textbf{lower Darboux sum} $L(f,P) = \sum_{i=1}^n \inf_{[x_{i-1},x_i]} f \cdot \Delta x_i$ under-estimates
the area, and the \textbf{upper Darboux sum} $U(f,P) = \sum_{i=1}^n \sup_{[x_{i-1},x_i]} f \cdot \Delta x_i$
over-estimates it.

Every lower sum is bounded above by every upper sum. Therefore, the set of all possible lower sums
$\set{L(f,P) : P \text{ a partition}}$ is nonempty and bounded above; completeness hands us its
supremum, the \textbf{lower integral} $I_*$. Symmetrically, the upper sums form a set bounded below;
completeness hands us its infimum, the \textbf{upper integral} $I^*$.

A bounded function is \emph{Riemann integrable} precisely when these two numbers meet: $I_* = I^*$.
For every continuous function on $[a,b]$, refining the partitions forces the gap $U(f,P) - L(f,P)$
to zero, pinching the lower and upper integrals together (Figure~\ref{ch41:fig-darboux}).
The common real value is the definite integral $\int_a^b f(x)\,\dd x$.
Without completeness, the set of lower sums would climb toward an empty hole, and the area under
the curve would have no real number to equal.

\begin{figure}[H]
\centering
\begin{tikzpicture}[scale=1.05]
  % Number line of area values
  \draw[->,thick] (0,0) -- (10.8,0) node[right] {Area};

  % True integral location at x = 5.4
  \def\intx{5.4}

  % Lower sums L(f, P) approaching from the left
  \foreach \x/\lab in {1.2/{L(f,P_1)}, 2.7/{L(f,P_2)}, 4.2/{L(f,P_3)}} {
    \fill[sourceblue] (\x, 0) circle (2.2pt);
    \draw[sourceblue, thick] (\x, 0.12) -- (\x, -0.12);
    \node[sourceblue, below=4pt, font=\footnotesize] at (\x, 0) {$\lab$};
  }
  \draw[sourceblue, thick, ->] (1.2, 0.45) -- (4.8, 0.45)
    node[midway, above, font=\small] {Lower sums climb: $\sup L(f,P)$};

  % Upper sums U(f, P) approaching from the right
  \foreach \x/\lab in {9.6/{U(f,P_1)}, 8.1/{U(f,P_2)}, 6.6/{U(f,P_3)}} {
    \fill[quotientorange] (\x, 0) circle (2.2pt);
    \draw[quotientorange, thick] (\x, 0.12) -- (\x, -0.12);
    \node[quotientorange, below=4pt, font=\footnotesize] at (\x, 0) {$\lab$};
  }
  \draw[quotientorange, thick, <-] (6.0, 0.45) -- (9.6, 0.45)
    node[midway, above, font=\small] {Upper sums descend: $\inf U(f,P)$};

  % Meeting point (The Integral)
  \draw[targetgreen!80!black, line width=1.5pt] (\intx, -0.7) -- (\intx, 0.9);
  \fill[targetgreen!80!black] (\intx, 0) circle (2.8pt);
  \node[targetgreen!80!black, above=2pt, font=\small\bfseries, align=center] at (\intx, 0.9)
    {The Integral $I = \int f(x)\,\dd x$};
  \node[targetgreen!80!black, below=5pt, font=\footnotesize, align=center] at (\intx, -0.7)
    {Meeting point: $I_* = I^*$\\Completeness guarantees this real value exists};
\end{tikzpicture}
\caption{Darboux integration via completeness. Every lower sum $L(f,P)$ (blue) underestimates
the area; every upper sum $U(f,P)$ (orange) overestimates it. As partitions refine, the lower sums
climb toward their supremum and the upper sums descend toward their infimum. For a continuous
function, these two values squeeze together to a single point. Completeness ensures that this
meeting point is an actual real number: the definite integral $\int_a^b f(x)\,\dd x$.}
\label{ch41:fig-darboux}
\end{figure}

\begin{aestheticbox}
Notice the common structure across all four theorems: the limit of a monotone sequence is a
supremum; the crossing point of a continuous curve is a supremum; the maximum of a continuous
function is an attained supremum; and the area under a curve is the common value where a
supremum of lower sums meets an infimum of upper sums. Each result applies the same
completeness property to a different mathematical setting. The completeness axiom provides
the unified foundation for analysis: every existence theorem for limits, derivatives, and
integrals relies on this single property of the real numbers.
\end{aestheticbox}

\section{The Archimedean property: no infinities hiding}
\label{ch41:archimedean}

Completeness implies another property that is used in nearly every $\eps$-argument: the
real numbers contain no infinitely large elements or positive infinitesimals.

\begin{theorem}[Archimedean property]
\label{ch41:arch}
For any real number $x$, there is a natural number $n$ with $n > x$. Equivalently, for any
$\eps>0$ there is a natural number $n$ with $\frac1n < \eps$.
\end{theorem}

\begin{proof}
Suppose not: suppose some real $x$ exceeds every natural number, so $\N$ is bounded above
by $x$. Then by completeness $\N$ has a least upper bound $L=\sup\N$. Since $L$ is the
\emph{least} upper bound, $L-1$ is not an upper bound, so some natural number $n$ satisfies
$n > L-1$. But then $n+1 > L$, and $n+1$ is also a natural number --- contradicting $L$
being an upper bound for $\N$. So no real number exceeds all naturals.

For the second form, given $\eps>0$ apply the first part to $x=1/\eps$ to get $n>1/\eps$,
whence $1/n<\eps$.
\end{proof}

\begin{happeningbox}
The Archimedean property guarantees that $\frac1n$ can be made smaller than any given positive
tolerance $\eps$ --- an essential step in $\eps$-$\delta$ arguments. It rules
out any positive real number that is strictly smaller than every $\frac1n$ (an infinitesimal),
because such a number would force $\N$ to be bounded. In this way, completeness ensures that
$\R$ contains no infinitely large elements or positive infinitesimals. (One
\emph{can} construct non-Archimedean ordered fields with genuine infinitesimals, such as the hyperreals
in nonstandard analysis, but they do not satisfy the completeness axiom.)
When we write ``for every $\eps>0$ there exists $N$,'' we rely on the Archimedean property,
which itself derives from completeness.
\end{happeningbox}

\begin{recoverybox}
\textbf{Pause and think:} Suppose someone claims to have found a positive real number $c > 0$ that is strictly smaller than $\frac{1}{n}$ for every single natural number $n \in \N$. Why does the Archimedean property immediately contradict this?
\emph{(Check: Apply Theorem~\ref{ch41:arch} to the real number $x = 1/c$. There must exist some $n \in \N$ with $n > 1/c$, which rearranges to $\frac{1}{n} < c$, immediately disproving the claim that $c$ is smaller than all fractions $\frac{1}{n}$.)}
\end{recoverybox}

\begin{connectionbox}
The Archimedean property also explains why the rationals, holes and all, still manage to be
\emph{dense} in the reals: between any two reals $a<b$ there is a rational. Choose $n$ with
$\frac1n < b-a$ (Archimedean), so the rationals $\frac{k}{n}$ are spaced more finely than
the gap $b-a$; one of them must land inside $(a,b)$. So every real number, $\sqrt2$ and
$\pi$ included, has rationals arbitrarily close to it on both sides. The irrationals fill
the holes, and the rationals crowd right up against every one of them. This is the precise
sense in which the decimal expansions of $\sqrt2$ and $\pi$ approximate them: rational
truncations creeping arbitrarily close to a real destination that completeness guarantees
is there.
\end{connectionbox}

\section{A worked tour of suprema and infima}
\label{ch41:worked}

Abstract definitions calcify unless you compute with them. Here are four sets; for each we
find $\sup$ and $\inf$ and decide whether they are attained. Work each one yourself before
reading the resolution.

\begin{example}[A set with inf attained, sup not]
\label{ch41:ex1}
Let $S = \set{1 - \tfrac1n : n\in\N} = \set{0,\ \tfrac12,\ \tfrac23,\ \tfrac34,\ \ldots}$.

\textbf{Infimum.} The smallest term is the first one, $1-\frac11 = 0$. Every term is
$\ge 0$, and $0$ is in $S$, so $\inf S = 0$ and it \emph{is attained} (it is the minimum).

\textbf{Supremum.} The terms climb: $0,0.5,0.667,0.75,0.8,\ldots$, always below $1$ since
$\frac1n>0$, yet creeping toward $1$ as $n$ grows (Archimedean: $\frac1n\to0$). So $1$ is
an upper bound, and it is the least one --- for any $\eps>0$, picking $n>\frac1\eps$ gives
a term $1-\frac1n > 1-\eps$, so no smaller ceiling works. Thus $\sup S = 1$, and it is
\emph{not attained}: no term equals $1$. Here is the open-interval phenomenon again, now in
a sequence.
\end{example}

\begin{example}[Sup of a set defined by an inequality]
\label{ch41:ex2}
Let $S = \set{x\in\R : x^2 < 5}$. This is the open interval $(-\sqrt5,\sqrt5)$. It is
bounded above by $\sqrt5\approx2.236$, and $\sup S=\sqrt5$ (not attained, since
$(\sqrt5)^2=5\not<5$). Symmetrically $\inf S = -\sqrt5$ (not attained). The supremum
$\sqrt5$ is itself an irrational number whose existence comes from completeness --- the same move
as Theorem~\ref{ch41:sqrt-exists} --- so this innocuous-looking set has a supremum that
does not exist at all in $\Q$.
\end{example}

\begin{example}[Attained extremes with unattained limit points]
\label{ch41:ex3}
Let $S = \set{(-1)^n\left(1+\tfrac1n\right) : n\in\N}$. The even terms are
$1+\frac1n$ (just above $1$, decreasing toward $1$): $\frac32, \frac54, \frac76,\ldots$
The odd terms are $-(1+\frac1n)$ (just below $-1$, increasing toward $-1$):
$-2, -\frac43, -\frac65,\ldots$

\textbf{Supremum.} The largest single term is the first even one, $n=2$: $1+\frac12=\frac32$.
All later even terms are smaller, all odd terms are negative. So $\sup S = \frac32$,
\emph{attained}.

\textbf{Infimum.} The smallest single term is the first odd one, $n=1$: $-2$. So
$\inf S = -2$, \emph{attained}. Both extremes here happen to be reached; the limit points
$\pm1$ that the sequence clusters around are \emph{not} the sup or inf --- a useful reminder
that ``where a sequence accumulates'' and ``its supremum'' are different questions.
\end{example}

\begin{example}[An unbounded set has no supremum]
\label{ch41:ex4}
Let $S = \set{n - \tfrac1n : n\in\N} = \set{0,\ \tfrac32,\ \tfrac83,\ \tfrac{15}{4},\ldots}$.
The terms grow without bound (by the Archimedean property, $n-\frac1n$ exceeds any fixed
number once $n$ is large). So $S$ is \emph{not bounded above} and has \emph{no supremum} ---
the completeness axiom never applies, since its hypothesis ``bounded above'' fails. Its
infimum is $\inf S=0$, attained at $n=1$. The lesson: completeness guarantees a supremum
\emph{only for bounded sets}; the hypothesis is not decorative.
\end{example}

\begin{labbox}
Try these before peeking at the answers in your head. For each set, find $\sup$ and $\inf$,
and say whether each is attained.
\begin{enumerate}
  \item $S = \set{\tfrac{n}{n+1} : n\in\N}$. \quad
        \emph{(Answer: $\inf=\tfrac12$ attained; $\sup=1$ not attained.)}
  \item $S = \set{x\in\Q : x^3 < 7}$. \quad
        \emph{(Answer: $\sup=\sqrt[3]{7}$, irrational, not attained; no inf, unbounded below.)}
  \item $S = \set{(-1)^n + \tfrac1n : n\in\N}$. \quad
        \emph{(Answer: $\sup=\tfrac32$ at $n=2$, attained; $\inf=-1$, the limit of odd
        terms, \emph{not} attained.)}
\end{enumerate}
Notice how often the supremum is the place the set is \emph{heading} rather than a place it
\emph{reaches}. That gap between heading-toward and reaching is the entire reason limits ---
and completeness --- exist.
\end{labbox}

\begin{applicationbox}
Numerical root-finding relies on this exact principle to compute $\sqrt2$. A numerical
root-finder never produces $\sqrt2$ exactly --- it produces a \emph{bounded, monotone
sequence} of rational approximations and trusts completeness to guarantee a real
destination. Newton's method (Exercise~\ref{ch41:ex-monotone}) generates
$1, 1.5, 1.41\overline{6}, 1.41421\ldots$, each a rational, each closer; the
\emph{number} $\sqrt2$ is the limit these approximations define. Every time your
calculator displays $1.41421356$, it is showing you a truncation of a real number whose
\emph{existence} is underwritten by the completeness axiom. Floating-point arithmetic lives
entirely among rationals; the meaning of its answers lives in $\R$.
\end{applicationbox}

\section{The shape of the whole idea}
\label{ch41:shape}

The logical skeleton is more memorable than any single proof, so gather the thread.

\begin{walkbox}
Walk the chain of dependence out loud:
\begin{enumerate}
  \item The rationals are dense but \emph{holey}: $\sqrt 2$ names a real length that no
        fraction reaches (Theorem~\ref{ch41:sqrt2-irrational}).
  \item A holey number system breaks limits: a bounded increasing sequence can head toward
        a destination the system does not contain.
  \item The cure is one axiom: \emph{every nonempty set bounded above has a least upper
        bound} (Section~\ref{ch41:axiom}). This declares the line hole-free.
  \item From that one axiom flow all the existence theorems: square roots
        (Theorem~\ref{ch41:sqrt-exists}), monotone convergence (Theorem~\ref{ch41:mct}),
        the IVT (Theorem~\ref{ch41:ivt}), the EVT (Theorem~\ref{ch41:evt}),
        Bolzano--Weierstrass (Theorem~\ref{ch41:bw}), the Archimedean property
        (Theorem~\ref{ch41:arch}), and the existence of the integral.
\end{enumerate}
Every ``there exists'' in this book traces back, through this chain, to step 3.
\end{walkbox}

\begin{aestheticbox}
Notice the historical reversal in how modern analysis is organized. Calculus came
first, and the rigorous construction of $\R$ came two centuries later --- engineered
specifically so that bounded monotone sequences and continuous curves crossing zero
would never fail for lack of points. When you open an analysis textbook, the
completeness axiom appears on page one as an abstract rule about sets and ceilings,
when historically it was reverse-engineered to make the theorems of calculus hold.
\end{aestheticbox}

\section{The Continuum: Triumphs, Costs, and the Horizon Ahead}
\label{ch41:what-next}

Completeness provides the rigorous foundation for calculus. By ensuring the absence of
gaps in the real line, the Least Upper Bound Property guarantees a limit for every bounded
monotone or Cauchy process. In $\Q$, a Cauchy sequence can fail to converge because its
apparent limit (such as $\sqrt{2}$ or $\pi$) is absent from $\Q$. Completeness guarantees
that every Cauchy sequence of real numbers converges to an element of $\R$, making $\R$ a
\textbf{complete metric space}.

The Intermediate Value Theorem, the Extreme Value Theorem, the Monotone Convergence Theorem,
and the Bolzano--Weierstrass Theorem are not isolated results; each of them is a direct
consequence of order completeness. The definite integral $\int_a^b f(x)\,\dd x$ likewise
relies on completeness: Riemann's construction requires the supremum of lower sums to equal
the infimum of upper sums, and completeness guarantees that these bounding numbers exist.
The foundational issue identified by the Pythagoreans---that geometric lengths cannot all be
expressed as rational ratios---is fully resolved: geometric continuity and numerical arithmetic
are unified.

Yet completeness carries philosophical and computational consequences.
The completeness axiom is non-constructive: it proves that $\sup S$ exists,
but it gives no mechanical formula to generate its digits. As Alan Turing proved in 1936,
the vast majority of real numbers are uncomputable---no finite algorithm running indefinitely
can compute their successive decimals. In addition, Georg Cantor's diagonal argument revealed
that $\R$ is \textbf{uncountably infinite} --- an infinity so large that its points cannot
be arranged into a numbered sequence $1, 2, 3, \ldots$ This led mathematics to foundational
questions such as the Continuum Hypothesis, which Kurt G\"odel and Paul Cohen proved to be
formally undecidable within standard ZFC set theory. Finally, while completeness provides
integrals for continuous functions, Riemann integration cannot integrate pointwise limits of
discontinuous functions, an obstacle that led Henri Lebesgue in 1902 to introduce
\textbf{measure theory} --- a general theory for assigning measure to sets by partitioning
the range rather than the domain.

\begin{connectionbox}
Calculus is the gateway to the rest of modern pure mathematics. With the real continuum
and the calculus of one and several variables firmly in hand, three natural pathways open
immediately ahead:
\begin{enumerate}
  \item \textbf{Linear Algebra From the Ground Up (Book 3):}
        In multivariable calculus (Chapters~32 and~33) we saw that derivatives of multivariable
        maps are linear transformations represented by Jacobian matrices. Linear algebra
        develops vector spaces, matrix factorizations, eigenvalues, and orthogonal projections
        in full generality, elevating the local linear approximations of calculus into a
        universal language for physics, optimization, and high-dimensional geometry.
  \item \textbf{Probability and Statistics From the Ground Up (Book 4):}
        Continuous probability density functions, the Gaussian distribution $\frac{1}{\sqrt{2\pi}}e^{-x^2/2}$,
        joint distributions, and the Central Limit Theorem build directly on the improper integrals,
        multiple integrals, and sequence limits established across this book.
  \item \textbf{Real Analysis and Measure Theory:}
        To push beyond Riemann's integral and complete function spaces into the \textbf{Hilbert space}
        $L^2$ (a space where functions are treated like infinite-dimensional vectors with well-defined
        lengths and angles), pure mathematics turns to Lebesgue integration, topology, and functional
        analysis---the mathematical language of modern quantum mechanics and signal analysis.
\end{enumerate}
\end{connectionbox}

\section{Exercises}
\label{ch41:exercises}

\begin{exercise}
\label{ch41:ex-finite-max}
\textbf{(Mental check: finite sets.)} Let $S = \{x_1, x_2, \ldots, x_k\}$ be any finite, nonempty set of real numbers.
Explain why $\sup S$ must not only exist, but must always be \emph{attained} (meaning $\sup S = \max S$).
Contrast this with infinite sets such as $(0,1)$ or $\{1 - 1/n : n \in \N\}$, where the supremum is never reached by any element of the set.
\end{exercise}

\begin{exercise}
\label{ch41:ex-sqrt3}
Adapt the proof of Theorem~\ref{ch41:sqrt2-irrational} to show $\sqrt 3$ is irrational.
(Hint: you need the fact that if $3 \mid p^2$ then $3 \mid p$. Show this by checking the
remainder of $p$ modulo $3$.) Then explain where the analogous argument \emph{breaks down}
for $\sqrt 4$, and why it must, since $\sqrt4=2$ is rational.
\end{exercise}

\begin{exercise}
\label{ch41:ex-sup-inf}
For each set, find $\sup$ and $\inf$ and state whether each is attained.
\[
\text{(a) } \set{2 + \tfrac{(-1)^n}{n} : n\in\N}, \quad
\text{(b) } (-3,\,5], \quad
\text{(c) } \set{x\in\R : x^2 - x < 0}, \quad
\text{(d) } \set{\tfrac{1}{n}+\tfrac{1}{m} : n,m\in\N}.
\]
\end{exercise}

\begin{exercise}
\label{ch41:ex-eps-char}
Use the $\eps$-characterisation of the supremum (Section~\ref{ch41:sup}) to prove the
following: if $S$ is nonempty and bounded above with $L=\sup S$, then there is a sequence
$(x_n)$ of elements of $S$ with $x_n \to L$. (Hint: for each $n$, apply tightness with
$\eps = 1/n$.)
\end{exercise}

\begin{exercise}
\label{ch41:ex-sup-sum}
Let $A$ and $B$ be nonempty sets bounded above, and define
$A+B = \set{a+b : a\in A,\ b\in B}$. Prove that $\sup(A+B) = \sup A + \sup B$. (Show
$\le$ using that $\sup A,\sup B$ are upper bounds, and $\ge$ using the $\eps$-tightness of
each, taking $\eps/2$ for each set.)
\end{exercise}

\begin{exercise}
\label{ch41:ex-ivt-rational}
Define $f:\Q\to\Q$ by $f(x)=x^2-2$. Verify that $f$ is continuous, $f(1)<0$, and $f(2)>0$,
yet $f$ has no zero in $\Q$. Explain in one or two sentences exactly which hypothesis of the
Intermediate Value Theorem (Theorem~\ref{ch41:ivt}) is unavailable here, and connect your
answer to the completeness axiom.
\end{exercise}

\begin{exercise}
\label{ch41:ex-monotone}
Define a sequence by $a_1 = 1$ and $a_{n+1} = \tfrac12\!\left(a_n + \tfrac{2}{a_n}\right)$
(this is Newton's method applied to $x^2-2=0$). Compute $a_1,a_2,a_3,a_4$ as fractions and
as decimals. Argue that for $n\ge 2$ the sequence is decreasing and bounded below by
$\sqrt2$, hence convergent by the Monotone Convergence Theorem
(Theorem~\ref{ch41:mct}), and show its limit $L$ must satisfy $L=\tfrac12(L+\tfrac2L)$, so
$L^2=2$. You have just built $\sqrt2$ from rationals using completeness.
\end{exercise}

\begin{exercise}
\label{ch41:ex-archimedean}
Use the Archimedean property (Theorem~\ref{ch41:arch}) to prove that
$\inf\set{\tfrac1n : n\in\N} = 0$. Be explicit about both promises an infimum must
satisfy: that $0$ is a lower bound, and that no positive number is a lower bound.
\end{exercise}

\begin{exercise}
\label{ch41:ex-nested}
\textbf{(Nested intervals.)} Suppose $[a_1,b_1] \supseteq [a_2,b_2] \supseteq \cdots$ is a
shrinking chain of closed intervals with lengths $b_n - a_n \to 0$. Use the completeness
axiom directly (consider $\sup\set{a_n : n\in\N}$) to prove there is exactly one real number
belonging to \emph{all} the intervals. Then explain why this fails over $\Q$ by exhibiting a
nested chain of rational intervals whose only common point would be $\sqrt2$.
\end{exercise}

\begin{exercise}
\label{ch41:ex-dense}
Prove that between any two real numbers $a < b$ there is an irrational number. (Hint:
between $a-\sqrt2$ and $b-\sqrt2$ there is a rational $r$ by density; consider $r+\sqrt2$.)
Combined with the density of the rationals, conclude that both the rationals and the
irrationals are dense in $\R$ --- every interval, however tiny, contains infinitely many of
each.
\end{exercise}

\begin{exercise}
\label{ch41:ex-cauchy-bounded}
\textbf{(The Cauchy criterion.)} A sequence $(x_n)$ is called a \emph{Cauchy sequence} if for every $\eps > 0$ there exists an integer $N$ such that $\abs{x_n - x_m} < \eps$ whenever $n, m \ge N$.
\begin{enumerate}
  \item Prove that every Cauchy sequence is bounded. (Hint: take $\eps = 1$, find the corresponding index $N$, and enclose all terms inside the bounds formed by the finite set $\set{x_1, \ldots, x_{N-1}}$ and the interval $(x_N - 1, x_N + 1)$.)
  \item Use the Bolzano--Weierstrass Theorem (Theorem~\ref{ch41:bw}) to extract a convergent subsequence $x_{n_k} \to L$, and prove that the full sequence $(x_n)$ converges to $L$.
\end{enumerate}
This establishes Cantor's theorem: in $\R$, a sequence converges if and only if its terms crowd arbitrarily close together.
\end{exercise}
